% Report body for the Hinglish sentiment project (SemEval-2020 Task 9, Hinglish).
% HOW TO USE: paste everything below into the body of the course's Overleaf template (between \begin{document}
% and \end{document}), after the title, authors and \maketitle. Upload the PNG files named in \includegraphics
% to Overleaf (they are in results/figures/ and in results/ of the repository). graphicx must be loaded.
% Text in \textbf{[PENDING ...]} marks work that has no result yet; replace or delete it before submission.
% Every number comes from a file in the repository; the source is named in a comment next to it.

\begin{abstract}
We study sentiment classification of Hindi--English (Hinglish) code-mixed tweets on the SemEval-2020 Task~9 data
(14{,}000 training, 3{,}000 development and 3{,}000 test tweets; negative, neutral, positive). We compare a
TF-IDF logistic-regression baseline with a BiLSTM over word and character embeddings, and ask two questions: does
normalising Romanised spelling help, and does the amount of code-mixing, measured by the Code-Mixing Index (CMI),
change accuracy? On the development set both model families reach about 0.61 weighted F1. The BiLSTM reaches
0.657--0.668 weighted F1 on the test set, evaluated once under a protocol fixed in advance. Spelling normalisation did not
improve results, and we find no clear evidence that CMI bucket is related to performance. Most errors involve the
neutral class, and about a quarter of a hand-read sample of errors have gold labels that look questionable.
\textbf{[PENDING: add one sentence on the transformer and LLM results if they are completed.]}
\end{abstract}

\section{Introduction}
Social-media text often mixes languages inside one sentence. Hinglish, Hindi written mostly in Roman script
and mixed with English, adds spelling variation on top of code-mixing: the same Hindi word appears as
\emph{acha}, \emph{accha} or \emph{achha}. SemEval-2020 Task~9 (SentiMix) asks systems to label such tweets as
negative, neutral or positive \cite{patwa2020}. The best system in the task reached 0.750 weighted F1, and annotators
agreed on only 55\% of tweets \cite{patwa2020}, so the task is hard and its labels are noisy.

Our project compares models of increasing strength and asks two research questions:
\begin{enumerate}
\item Does normalising Romanised spelling (transliteration normalisation) change accuracy?
\item Does the amount of code-mixing in a tweet change accuracy?
\end{enumerate}
\textbf{Contributions.} (i) A reproducible pipeline with fixed splits, a single metric implementation, tests and a
frozen one-time test protocol. (ii) A comparison of TF-IDF and BiLSTM models, with three preprocessing versions of the
text and three random seeds each. (iii) An analysis by code-mixing bucket. (iv) A categorised error analysis.
\textbf{[PENDING: add the fine-tuned transformers (mBERT, XLM-R, MuRIL/IndicBERT) and zero-shot LLM prompting if
completed; both are part of the project plan but had no results when this draft was written.]}

\section{Related work}
\textbf{Task and data.} SentiMix \cite{patwa2020} released code-mixed tweets labelled by sentiment and with
per-word language tags. \textbf{Code-mixing measures.} The Code-Mixing Index of Gamb\"ack and Das \cite{gamback2014}
measures how evenly languages are mixed in an utterance. \textbf{Pretrained multilingual models.} mBERT
\cite{devlin2019} and XLM-R \cite{conneau2020} are multilingual encoders; MuRIL \cite{khanuja2021} and IndicBERT
\cite{kakwani2020} are trained for Indian languages.
\textbf{[PENDING: the team's paper summaries (docs/paper\_summary\_template.md) have not been filled in. Write 2--3
paragraphs on SentiMix systems and on code-mixed sentiment work from them.]}

\section{Dataset and analysis}
\subsection{Data}
We use the Hinglish subset of SemEval-2020 Task~9 \cite{patwa2020} (Zenodo record 3974927, CC-BY 4.0), included in the
repository. Each tweet has one label and per-word language tags (Hindi, English, other, emoticon). The text is tokenised
and mostly in Roman script; only 59 of 14{,}000 training tweets contain Devanagari. The test labels are in a separate
file and the test set is used only for final evaluation.

\subsection{Labels and metric}
Classes are close to balanced (Table~\ref{tab:labels}). A classifier that always predicts neutral reaches 37.6\%
accuracy but only 0.205 weighted F1 on development data, which is the floor. We report the official metric of the
task, \textbf{weighted F1} \cite{patwa2020}, together with macro F1 and accuracy.

\begin{table}[h]
\centering
\begin{tabular}{lrrrr}
\hline
Split & Tweets & Negative & Neutral & Positive \\
\hline
Train & 14{,}000 & 29.3\% & 37.6\% & 33.1\% \\
Dev   & 3{,}000  & 29.7\% & 37.6\% & 32.7\% \\
Test  & 3{,}000  & 30.0\% & 36.7\% & 33.3\% \\
\hline
\end{tabular}
\caption{Label distribution (docs/eda\_part1.md).}
\label{tab:labels}
\end{table}

\subsection{Text properties}
Tweets are short (mean 26 word tokens, maximum 56). Hindi tokens make up 46.5\% of all tokens, English 33.2\%
and other tokens (mentions, URLs, punctuation) 20.2\%. In the training data 56\% of tweets contain a URL fragment, 81\% a
mention, 16\% start with ``RT'' and 11\% are cut off at the end. Eleven training tweets are almost entirely in one
language, and no train--dev or dev--test text is shared, so leakage between splits is negligible.

\subsection{Code-Mixing Index}
We compute $\mathrm{CMI}=100\,(1-\max(w)/(n-u))$ per tweet, where $n$ is the number of tokens, $u$ the number of
language-independent tokens and $\max(w)$ the count of the most frequent language \cite{gamback2014}. With two languages the
maximum is 50. Our mean on the training set is 29.8, higher than the 25.3 reported by the task organisers
\cite{patwa2020}. We tried several variants of the formula and none reproduced their value, so we use the standard
formula consistently and compare tweets by rank only. Almost every tweet is mixed (the 10th percentile is 16.7), so our
``low'', ``medium'' and ``high'' groups, the training-set terciles with cut points 23.8 and 34.8, compare
moderately mixed with heavily mixed text and not monolingual with mixed text (Figure~\ref{fig:cmi}).

\begin{figure}[h]
\centering
\includegraphics[width=0.7\linewidth]{cmi_distribution.png}
\caption{Distribution of CMI on the training set (results/figures/cmi\_distribution.png).}
\label{fig:cmi}
\end{figure}

\section{Methodology}
\subsection{Experimental protocol}
Models are trained on the training set. The development set is used for early stopping and for comparing settings. The
test set is used once, for final evaluation, under a protocol written before the test labels were used for any
model purpose; test predictions are saved before test labels are read, and no result selects a model, seed or setting.
Each BiLSTM configuration was trained with seeds 42, 43 and 44. We report mean $\pm$ sample standard deviation over seeds.
The unit tests and the team's exploratory data analysis loaded the test split for descriptive statistics, and this is
disclosed in the repository (docs/bilstm\_final\_evaluation\_protocol.md).

\subsection{Preprocessing versions}
We compare three versions of the same tweets. \emph{Raw}: text as released. \emph{Clean}: URLs and mentions replaced by
the placeholders URL and USER, hashtags kept without the \#, retweet markers and truncation marks removed, split
contractions rejoined. \emph{Normalized}: clean, plus runs of the same letter shortened to two and a dictionary that maps
spelling variants of common Romanised Hindi words to one form (25 canonical words, 57 variants). The dictionary changes
14.6\% of training tweets beyond cleaning. It excludes variants that are also common English words.

\subsection{TF-IDF baseline}
Word 1--2-gram, character 2--5-gram and combined TF-IDF features with logistic regression. Vocabulary and IDF weights are
fitted on the training set only, and the regularisation strength was chosen on the development set.

\subsection{BiLSTM}
Each word is represented by a 100-dimensional learned word embedding concatenated with a character CNN (kernel widths
2, 3, 4; 32 filters each) over its letters, so spelling variants and unseen words still receive a useful vector. A one-layer
BiLSTM (128 units per direction) reads the tweet, and max and mean pooling over words feeds a linear layer with three outputs.
Dropout is 0.4; training uses Adam (learning rate 0.001), batch size 64, at most 20 epochs and early stopping with patience~4
on development weighted F1. Embeddings are trained from scratch. Vocabularies and length limits (56 words, 18 characters) are
derived from the training set only and are shared by the three preprocessing versions.

\subsection{Transformers and LLM prompting}
\textbf{[PENDING: describe the fine-tuned mBERT, XLM-R and MuRIL/IndicBERT models (shared training pipeline, learning rates,
epochs, seeds, maximum subword length) and the zero-shot LLM prompt (model, prompt text, output parsing) once they exist.
No results were available when this draft was written.]}

\section{Experiments and results}
\subsection{Development set}
Table~\ref{tab:dev} shows development results. All models are within about 0.006 weighted F1 of each other. The typical
sampling error of an F1 score on 3{,}000 tweets is about 0.009, so we cannot separate the model families, the TF-IDF feature
sets or the three BiLSTM text versions on this set.

\begin{table}[h]
\centering
\begin{tabular}{lrrr}
\hline
Model & Weighted F1 & Macro F1 & Accuracy \\
\hline
Always ``neutral'' & 0.205 & 0.182 & 0.376 \\
TF-IDF word 1--2 & 0.611 & 0.616 & 0.612 \\
TF-IDF char 2--5 & 0.613 & 0.616 & 0.612 \\
TF-IDF word + char & 0.612 & 0.615 & 0.611 \\
BiLSTM raw (3 seeds) & 0.611 $\pm$ 0.004 & 0.614 $\pm$ 0.006 & 0.610 $\pm$ 0.005 \\
BiLSTM clean (3 seeds) & 0.612 $\pm$ 0.002 & 0.617 $\pm$ 0.002 & 0.613 $\pm$ 0.003 \\
BiLSTM normalized (3 seeds) & 0.607 $\pm$ 0.001 & 0.611 $\pm$ 0.000 & 0.606 $\pm$ 0.001 \\
\hline
\end{tabular}
\caption{Development-set results (results/baseline\_scores.csv, results/experiment\_log.csv, reports/bilstm\_final\_analysis.md).
BiLSTM dev scores are the best epoch per run and are optimistically selected.}
\label{tab:dev}
\end{table}

\subsection{Test set}
Table~\ref{tab:test} gives the single frozen test evaluation of the nine BiLSTM runs. \textbf{TF-IDF was not evaluated on
the test set}, so there is no test comparison between model families.
\textbf{[PENDING: add rows for the transformer and LLM systems.]}

\begin{table}[h]
\centering
\begin{tabular}{lrrr}
\hline
BiLSTM text version & Weighted F1 & Macro F1 & Accuracy \\
\hline
Raw        & 0.660 $\pm$ 0.007 & 0.662 $\pm$ 0.008 & 0.658 $\pm$ 0.005 \\
Clean      & 0.668 $\pm$ 0.002 & 0.672 $\pm$ 0.003 & 0.667 $\pm$ 0.004 \\
Normalized & 0.657 $\pm$ 0.006 & 0.659 $\pm$ 0.006 & 0.653 $\pm$ 0.007 \\
\hline
\end{tabular}
\caption{Test-set results, mean $\pm$ sample standard deviation over seeds 42--44 (results/bilstm\_final\_test/). For reference,
the best SemEval Hinglish system reached 0.750 weighted F1 \cite{patwa2020}.}
\label{tab:test}
\end{table}

Test scores are 0.047 to 0.054 macro F1 higher than development scores for every text version. We checked uid alignment,
label mapping, class mix, tweet length, code-mixing and out-of-vocabulary rates, and none explains the gap; re-weighting the
test set to the development class mix changes macro F1 by less than 0.001. Whether the test tweets are simply easier, or the
gap is specific to this model family, is untested.

\section{Ablations}
\subsection{Spelling normalisation}
Normalisation did not help (Tables~\ref{tab:dev} and~\ref{tab:test}). On the test set a paired bootstrap over tweets gives
a difference in macro F1 of $-0.003$ (95\% interval $-0.012$ to $+0.006$) between normalized and raw text and
$-0.013$ ($-0.021$ to $-0.004$) between normalized and clean text; cleaning alone was the best version, by $+0.010$
($+0.0003$ to $+0.019$) over raw text in macro F1, but this interval just excludes zero and the weighted-F1 interval
contains it. The intervals describe test-sampling uncertainty only; seed-to-seed differences are of similar size and we
do not claim statistical significance. The dictionary is small (25 words), so this result says that this small dictionary
did not help, and not that spelling normalisation cannot help.

\subsection{Code-mixing}
We scored the development predictions of every model in three CMI buckets (about 1{,}000 tweets each) and estimated, with a
bootstrap, the difference in weighted F1 between buckets. Of 36 pairwise intervals (12 prediction files $\times$ 3
comparisons), 35 contain zero; one does not (BiLSTM on normalized text, seed 43, low versus high), which is about the
number expected by chance among 36 intervals at 95\%. High CMI is nominally the best bucket for all three TF-IDF models,
but the gaps (0.01 to 0.03) are smaller than the sampling error of a bucket (about 0.03). We therefore find \emph{no
clear evidence} of a relationship between CMI bucket and performance. With 1{,}000 tweets per bucket only differences of
roughly 0.04 to 0.05 or more could be detected.

\section{Error analysis}
\subsection{Where the errors are}
For the TF-IDF model with character features, 88\% of development errors involve the neutral class (1{,}021 of 1{,}163);
only 12\% swap negative and positive. The BiLSTM shows the same pattern: 89--91\% of development errors and 93--95\% of
test errors involve neutral. Neutral is the weakest class (TF-IDF F1 0.55 against 0.63 negative and 0.67 positive).

\subsection{What goes with errors}
Tweets containing a negation word are misclassified more often (44.1\% against 36.1\%; $p<0.001$, below the Bonferroni
threshold of $0.05/6$). Spelling variants from the normalisation dictionary, URLs, truncation and repeated letters show no
reliable link to errors. The negation flag is a keyword heuristic.

\subsection{Categories read by hand}
We read a random sample of 60 development errors of the TF-IDF model and gave each one main category
(Table~\ref{tab:cats}). This is a first-pass reading by one reader, and with 60 tweets each share is uncertain by about
$\pm10$ points. \textbf{[PENDING: add raw agreement and Cohen's kappa once a second reader has labelled the same 60 tweets
(python -m src.label\_agreement).]}

\begin{table}[h]
\centering
\begin{tabular}{lrr}
\hline
Category & Tweets & Share \\
\hline
Label looks questionable & 14 & 23\% \\
Context-dependent (no explicit sentiment words) & 13 & 22\% \\
Sarcasm or irony & 9 & 15\% \\
Negation & 8 & 13\% \\
Slang and informal language & 7 & 12\% \\
Transliteration variants & 5 & 8\% \\
Clear sentiment missed & 3 & 5\% \\
Code-mixing (mixed script) & 1 & 2\% \\
\hline
\end{tabular}
\caption{Hand-read categories for 60 TF-IDF development errors (results/error\_categories\_manual.csv).}
\label{tab:cats}
\end{table}

About a quarter of the sampled errors have a gold label that looks wrong, for example greetings and praise labelled
neutral. Together with the low annotator agreement \cite{patwa2020} this suggests a noise ceiling on every model; it is a
hypothesis until a second reader agrees. The most confident TF-IDF mistakes point the same way: 26 of the 30 are tweets
labelled neutral that read as clearly positive or negative. Only 1 of 60 errors falls in the code-mixing category, which fits the flat CMI result.

\section{Conclusion and limitations}
On Hinglish SentiMix, a BiLSTM with word and character embeddings reaches about 0.66 weighted F1 on the test set, and on
development data it is indistinguishable from a TF-IDF baseline. A small spelling-normalisation dictionary did not help,
and we find no clear evidence that code-mixing level changes accuracy. Errors concentrate at the boundary of the neutral
class, and part of that boundary may reflect label noise.

\textbf{Limitations.} Only the BiLSTM was evaluated on the test set. Three seeds per configuration give rough standard
deviations. Our CMI values differ from the organisers' and every tweet is mixed, so the CMI groups are not monolingual
versus mixed. The hand-read categories come from one reader. The dev-to-test gap is unexplained. The best SemEval system
(0.750 weighted F1) is well above our scores.
\textbf{[PENDING: state what was and was not done for the transformers and the LLM.]}

\begin{thebibliography}{9}
\bibitem{patwa2020} Patwa, P., Aguilar, G., Kar, S., Pandey, S., PYKL, S., Gamb\"ack, B., Chakraborty, T., Solorio, T.,
and Das, A. (2020). SemEval-2020 Task 9: Overview of Sentiment Analysis of Code-Mixed Tweets. In \emph{Proceedings of
the 14th International Workshop on Semantic Evaluation (SemEval-2020)}. arXiv:2008.04277.
\bibitem{gamback2014} Gamb\"ack, B. and Das, A. (2014). On Measuring the Complexity of Code-Mixing. In
\emph{Proceedings of the 11th International Conference on Natural Language Processing, Goa, India}.
\bibitem{devlin2019} Devlin, J., Chang, M.-W., Lee, K., and Toutanova, K. (2019). BERT: Pre-training of Deep Bidirectional
Transformers for Language Understanding. In \emph{Proceedings of NAACL-HLT 2019}.
\bibitem{conneau2020} Conneau, A., et al. (2020). Unsupervised Cross-lingual Representation Learning at Scale. In
\emph{Proceedings of ACL 2020}. arXiv:1911.02116.
\bibitem{khanuja2021} Khanuja, S., et al. (2021). MuRIL: Multilingual Representations for Indian Languages.
arXiv:2103.10730.
\bibitem{kakwani2020} Kakwani, D., et al. (2020). IndicNLPSuite: Monolingual Corpora, Evaluation Benchmarks and
Pre-trained Multilingual Language Models for Indian Languages. In \emph{Findings of EMNLP 2020}.
% VERIFY every reference above (authors, venue, year, pages) against the papers before submission; they were written
% from memory and not checked online.
\end{thebibliography}
